\subsection{Critical examples with no positive-dimensional regular subset}
\label{criticalmoran:section}

The covering gain in Theorem~\ref{endpoint:subpower} can occur at exact
critical dimensions. The following construction also prevents reducing
these examples to the previous strict condition by selecting a subset.

\begin{theorem}\label{criticalmoran:theorem}
For every \(1<d\le3/2\), put \(D=2d-1\). There are a compact
\(K\subset[0,1]^2\) and a \(d\)-Frostman probability \(\mu\) on \(K\)
such that
\begin{equation}\label{criticalmoran:dimensions}
 0<\mathcal H^d(K)<\infty,\qquad \HH K=d,\qquad \PP K=D,
\end{equation}
and, for all sufficiently small \(r>0\),
\begin{equation}\label{criticalmoran:cover}
 N(K,r)\le Cr^{-D}\exp\!\left(-c\sqrt{\log(1/r)}\right).
\end{equation}
Every subset \(S\subset K\) satisfies
\begin{equation}\label{criticalmoran:hereditary}
 \HH S\le\frac dD\,\PP S.
\end{equation}
In particular \(K\) has no subset of positive equal Hausdorff and
packing dimensions. Nevertheless its raw joint distance law with any
compactly supported pin Frostman probability of exponent greater than
one is absolutely continuous. In particular,
\[
 |\Delta_y(K)|>0\qquad\text{for \(\mu\)-almost every }y\in K.
\]
If \(d\le d_\alpha=(7-\sqrt7)/4\), no subset \(S\subset K\) with
\(\HH S>1\) satisfies the previous strict inequality
\(\PP S<B_H(\HH S)\).
\end{theorem}

\begin{proof}
Fix \(A>0\) and set \(U(n)=Dn-A\sqrt{n+1}\). Choose a large
integer \(a_0\) so that, for all \(n\ge a_0\),
\[
 U(n)-dn>10,\qquad 0<U(n+1)-U(n)<2.
\]
This is possible because \(D>d\) and \(D\le2\). The strict upper
increment bound holds even when \(D=2\), due to the square-root term.
Set \(F(a_0)=2\lceil da_0/2\rceil\), choosing \(a_0\) large enough
that \(F(a_0)\le2a_0\), and join \(F(0)=0\) to this even integer
using increments in \(\{0,2\}\).

Starting at \(a_{j-1}\), let \(F\) increase with slope two at integer
depths until its first crossing of \(U\). Explicitly, \(b_j\) is
the first integer \(n>a_{j-1}\) for which
\[
 F(a_{j-1})+2(n-a_{j-1})\ge U(n).
\]
This crossing is finite: for \(D<2\) the difference tends to
infinity linearly, and for \(D=2\) it is a constant plus
\(A\sqrt{n+1}\). The first-crossing property gives
\begin{equation}\label{criticalmoran:upperhit}
 U(b_j)\le F(b_j)<U(b_j)+2.
\end{equation}
Next keep \(F\) constant until the lower line first catches it:
let \(a_j>b_j\) be the first integer with \(da_j\ge F(b_j)\).
Then
\begin{equation}\label{criticalmoran:lowerhit}
 F(a_j)=F(b_j),\qquad da_j-d<F(a_j)\le da_j.
\end{equation}
The separation of the barriers makes both phases nonempty. Thus
\(a_j,b_j\to\infty\).

Every increment of \(F\) is zero or two, and for \(n\ge a_0\),
\begin{equation}\label{criticalmoran:barriers}
 dn-2\le F(n)\le Dn-A\sqrt{n+1}+2.
\end{equation}
For the lower bound the growing slope exceeds \(d\), while the
flat phase stops within one step of the lower line. For the upper
bound use \eqref{criticalmoran:upperhit} and the monotonicity of \(U\).
The endpoint identities also give
\begin{equation}\label{criticalmoran:ratio}
 da_j=Db_j-A\sqrt{b_j+1}+O(1),\qquad
 \frac{a_j}{b_j}\longrightarrow\frac Dd.
\end{equation}

At each planar dyadic depth \(n\), use the same rule in every
occupied square: retain all four children if
\(F(n)-F(n-1)=2\), and only the lower-left child otherwise.
Give equal conditional weights to the retained children.
The closed nested construction defines a compact \(K\) and a
probability \(\mu\). There are \(2^{F(n)}\) symbolic depth-\(n\)
cells, each of mass \(2^{-F(n)}\).

There are infinitely many zero-increment depths after every fixed
depth. Consequently no admissible coordinate tail can consist
entirely of ones. The alternative upper-endpoint expansion of a
dyadic rational is therefore inadmissible, and the half-open dyadic
partition gives unambiguous addresses with precisely the stated
cell masses. Equivalently, descendants in any fixed construction
square stay strictly below its upper edges because of a later
common zero digit. This justifies the use of exact cell counts
without discarding boundary points.

A ball of radius \(2^{-n}\) meets a bounded number of dyadic
depth-\(n\) cells. Equation~\eqref{criticalmoran:barriers} gives
\[
 \mu(B(x,2^{-n}))\le C2^{-F(n)}\le C2^{-dn}.
\]
The finite initial range changes only \(C\), and adjacent dyadic
scales prove the \(d\)-Frostman bound at arbitrary radii.
Similarly,
\[
 N(K,2^{-n})\le C2^{Dn-A\sqrt{n+1}},
\]
which gives \eqref{criticalmoran:cover}, allowing a fixed depth
shift to account for square diameter.

The mass distribution principle gives \(\mathcal H^d(K)>0\).
At the flat endpoints, the full cover by construction squares has
\(d\)-dimensional cost
\[
 2^{F(a_j)}(\sqrt2\,2^{-a_j})^d
 =2^{d/2}2^{F(a_j)-da_j}\le2^{d/2}.
\]
Its diameters tend to zero, so \(\mathcal H^d(K)<\infty\) and
\(\HH K=d\).

The upper covering estimate gives \(\PP K\le D\). Conversely,
at the growing endpoints,
\[
 F(b_j)=Db_j-A\sqrt{b_j+1}+O(1),\qquad
 \sup_x\mu(B(x,2^{-b_j}))\le C2^{-F(b_j)}.
\]
Every subset \(E\subset K\) of positive \(\mu\) outer measure
therefore obeys
\[
 N(E,2^{-b_j})\ge c\,\mu^*(E)\,2^{F(b_j)},
\]
so its upper box dimension is at least \(D\). Every countable
cover of \(K\) has a member of positive \(\mu\) outer measure.
The countable-cover characterization of packing dimension gives
\(\PP K\ge D\), proving \eqref{criticalmoran:dimensions}.

To prove the stronger hereditary statement, take any nonempty
\(S\subset K\) and write \(v=\PP S\). For \(\epsilon>0\), choose
a countable cover by \(S_i\subset K\) with upper box dimensions
at most \(v+\epsilon\). For every fixed \(i\), and all sufficiently
large \(j\),
\[
 N(S_i,2^{-b_j})\le2^{(v+2\epsilon)b_j}.
\]
Replacing the covering balls by the bounded number of dyadic
depth-\(b_j\) squares they meet changes only a constant. There is
no branching between \(b_j\) and \(a_j\). Every occupied
depth-\(b_j\) square has exactly one occupied depth-\(a_j\)
descendant, so \(S_i\) can be covered by at most
\(C2^{(v+2\epsilon)b_j}\) squares of side \(2^{-a_j}\).
For \(t>(d/D)(v+2\epsilon)\), equation~\eqref{criticalmoran:ratio}
shows that their \(t\)-dimensional cost tends to zero:
\[
 C2^{(v+2\epsilon)b_j}(\sqrt2\,2^{-a_j})^t\longrightarrow0.
\]
Thus \(\HH S_i\le(d/D)(v+2\epsilon)\). Countable stability of
Hausdorff dimension, followed by \(\epsilon\downarrow0\), proves
\eqref{criticalmoran:hereditary}. The empty case is immediate.
Since \(d/D<1\), this rules out positive-dimensional subsets with
equal Hausdorff and packing dimensions.

Theorem~\ref{endpoint:subpower} now applies to the natural
\(d\)-Frostman measure and its covering estimate. It proves the
asserted joint absolute continuity and the self-pinned
positive-length conclusion.

Finally suppose \(d\le d_\alpha\), and take \(S\subset K\) with
\(h=\HH S>1\). Then \(h\le d\le d_\alpha\), so
\(B_H(h)=2h-1\). The hereditary inequality gives
\[
 \PP S\ge\frac Dd\,h,\qquad
 \frac Dd\,h-(2h-1)=1-\frac hd\ge0.
\]
Hence no such subset satisfies the preceding strict dimension
criterion.
\end{proof}

These are additional endpoint examples under the quantitative covering
gain \eqref{criticalmoran:cover}. The equality \(\PP K=2\HH K-1\)
alone is not asserted to imply positive pinned length.
